\section{Introduction}\label{sec:intro}

\subsection{The problem}

Decision makers rarely receive unlimited opportunities to revise a plan.
A commander who can re-issue manoeuvre orders at every moment is not the same
agent as one who may re-issue them twice in an engagement, even when both
retain the same ships, the same sensors and the same doctrine. What differs is
the \emph{revision calendar}: the set of epochs at which the agent may look at
what has actually happened and choose again. Replanning is not free, it cannot
be scheduled after the fact, and the value of a revision depends on what the
opponent has already done.

This paper studies the resulting optimisation problem in a class of
adversarial trajectory games. Two formations move under nonholonomic
constraints inside a directional interaction field, and each earns a payoff
that integrates the field along the whole trajectory rather than at a terminal
instant. The state that matters is a history: trailing vessels follow the
leader's executed track, so the current leader pose does not determine the
formation's shape \citep{stengel1996}. Each player is allowed a bounded number
of revisions, arranged on a calendar that is fixed and public in advance.
Within each interval between consecutive revisions the player seals a block of
commands; the opponent sees the block only as it is executed.

The question is deliberately budget-shaped. Given a tolerated loss of
performance, how many revisions are needed, when should they occur, and can a
computation certify that the chosen calendar is good enough? Three things make
this non-trivial. First, the feasible policy set depends on the calendar, so
calendars cannot be compared by evaluating a single pre-computed policy.
Second, more revisions need not have diminishing returns: the natural greedy
argument fails, because revision opportunities can be complementary. Third,
the opponent may itself be revising, and a strategy that is optimal against a
committed opponent is not admissible against a flexible one.

\subsection{Approach}

We treat the calendar as an explicit decision variable and build an exact
finite-game formulation around it. Section~\ref{sec:model} fixes the
history-dependent interaction game, the sealed-block information structure and
the budgeted value $W_K$. Section~\ref{sec:theory} develops four structural
results and the certificates the computation relies on:
policy inclusion gives monotonicity in the budget and identifies the endpoints;
a counterexample shows that the calendar objective is not submodular, so
selection needs more than a greedy argument; a common-update decomposition
localises the game at public boundaries; and two certificate constructions
bound the loss of a coarsened calendar from above, one by an explicit feasible
policy and one by summing losses over public intervals. Section
\ref{sec:algorithm} turns these into a solver: exact pricing of revisions
against a fixed causal opponent, calendar branch and bound with certified
intervals, and reuse of opponent responses across calendars.
Section~\ref{sec:methods} freezes the experimental design, and
Section~\ref{sec:results} reports it.

\subsection{Contributions and scope}

\begin{enumerate}
\item \textbf{A certified budget frontier for history-dependent sealed-block
games.} For a fixed opponent class we compute, exactly up to a certified
interval, the value attainable with at most $K$ revisions, $W_K$, over a
frozen finite design. The frontier identifies which revisions actually pay and
where the budget stops binding.
\item \textbf{Structure that survives contact with counterexamples.} We prove
budget monotonicity and endpoint equivalence, exhibit a binary-action
counterexample showing the calendar objective is not submodular, and prove a
common-update decomposition whose scope is stated exactly: a revision by one
player alone is not a public boundary while the opponent still holds an
unexecuted block.
\item \textbf{Two computable loss certificates and an exact pricing
recurrence.} A constructive deletion certificate gives a feasible coarse
policy and an upper bound on the loss it incurs; an interval-loss bound gives
a shortest-path certificate for a budget under a fully flexible opponent. Both
are reported with the residual of the computation that produced them.
\end{enumerate}

The scope is bounded and we state it where it binds. The physical design is a
\emph{fixed computational benchmark} over three geometries, four
speed/distance configurations, two opponent classes and three one-factor
ablations; it is not a statistical sample and we do not attach $p$-values to
deterministic solver output. The opponent classes are a committed class and a
fully flexible class, and the second certificate applies only to the latter,
where public boundaries exist. A previously investigated claim --- that the
range-advantaged side values revision rights differently from the
range-disadvantaged side --- was tested at this fidelity and \emph{falsified};
those results stay archived and are disclosed as a negative finding rather
than repaired. Our claim is about the frontier and the certificates, not about
a new general equilibrium-computation method: the pricing recurrence
specialises established extensive-form best-response ideas
\citep{mcmahan2003,bosansky2014,kroer2018}, and completeness of the reported
frontier is relative to the frozen design and the certified tolerances.

\subsection{Related work}

Our setting sits between three literatures. In \emph{dynamic games with
information}, an informed player's value is monotone in the information
structure \citep{hogeboom2020}; revisions are not free observations but
binding command blocks, so the inclusion argument applies to policy sets
rather than to signal structures. Huang and Zhu study a pursuit--evasion
differential game in which observation opportunities are purchased
strategically \citep{huang2021}; there the cost is paid for information about
the state, whereas here the budget limits when the agent may \emph{act on}
information it already has. Chen and Waggoner show that whether additional
signals are substitutes or complements depends on the decision problem
\citep{chen2017}; our counterexample extends that caution to revision
calendars.

In \emph{zero-sum game solving}, double-oracle and column-generation methods
converge finitely with response oracles \citep{mcmahan2003,bosansky2014}, and
extensive-form abstractions admit bounds on equilibrium error
\citep{kroer2018}. We use those tools rather than reinvent them. The
distinction is the object being optimised: not a strategy within a fixed
information structure, but the information structure itself, subject to an
explicit revision budget and with per-calendar certificates.

In \emph{formation control and pursuit--evasion}, leader--follower
path-following formations and speed asymmetry under turn constraints are
classical. We use a leader--follower line-ahead formation as the numerical
instantiation and treat a rigid column as an ablation, so that conclusions are
not an artefact of one formation representation.
